% Section 5: districts and the rung ladder (PHASE2b)
\section{Districts and the rung ladder}\label{sec:districts}

Source: PHASE2b (Appendix~\ref{app:phase2b}). Checks:
\texttt{verify\_phase2b.py}, 22/22 PASS. ``District'' means any partition of the
units into groups (ridings, counties, schools, or election periods). The district
model replaces (S2) by $\E[\beta\mid V]=M_{H^{(D)}}(X)$ and (S3) by validity on
each district hull $A_d$. The rungs are restrictions on
$\mathbf H=(H^{(1)},\dots,H^{(D^\star)})$:
\begin{itemize}[nosep]
\item \textbf{M2}: none (each district its own linear effects).
\item \textbf{M1($\delta$)}: $\exists\bar k$ with $\|k^{(d)}-\bar k-c_d\|_\infty\le\delta$ for all $d$.
\item \textbf{M1}: $H^{(d)}_c=H_c+\zeta_{c,d}\mathbf 1^\top$, $\sum_c\zeta_{c,d}=0$, $\zeta_{c,1}=0$.
\item \textbf{L4($r$)}: M1 and row $r$ has no contextual effect.
\item \textbf{L5}: $H^{(d)}_c=\alpha_{c,d}\mathbf 1^\top$ (constancy within districts).
\item \textbf{L6}: L5 with common $\alpha$ (Goodman).
\item \textbf{M0}: one $H$ for all units on $\conv\bigcup_dA_d$ (not nested).
\end{itemize}
$\delta$ is in probability units: $\delta=0.05$ lets the gap between ``$r$-voters
in pure-$s$ polls'' and ``$s$-voters in pure-$r$ polls'' differ across districts
by up to 10 points.

\begin{result}[Proposition D0, grouping adds nothing without a model --- PHASE2b \S1]\label{res:d0}
\status{Proved.} Under (S1) alone the sharp set for $B^*_{rc}$ is the
Duncan--Davis interval $DD^*_{rc}$, whatever the law of $D$ or of any
covariate. The joint sharp set for the whole matrix $B^*$ is convex and compact
but not the product of the cell intervals.
\end{result}
\begin{plain}
District dummies cannot narrow the model-free bound: it is already computed unit
by unit. Every gain from districts comes from an assumption that links districts.
\end{plain}
\proofref{Appendix~\ref{app:phase2b}, \S1, steps 1--8 (Gale's theorem)}
\begin{revnote}(R-18.) The per-cell sharp interval is now separated from the joint set.\end{revnote}

\begin{result}[Lemma D1, what each rung identifies --- PHASE2b \S2]\label{res:d1}
\status{Proved.} Under A5 within districts: (a) every rung identifies
$Q^{(d)}$; each $K^{(d)}\in\Kspace$ is unidentified. (b) Under M1,
$\zeta_{c,d}=\operatorname{diag}(Q^{(d)}_c-Q^{(1)}_c)$ is identified; the off-diagonal entries
satisfy the overidentifying restrictions
$(Q^{(d)}_c-Q^{(1)}_c)_{rs}=\tfrac12(\zeta_{c,d,r}+\zeta_{c,d,s})$; the unknown
dimension is the same as pooled. (c) Under L5, $K^{(d)}$ is identified and the
structure of $Q^{(d)}$ is testable. (d) Under M1 a weaker stacked rank condition
(A5-M1) suffices; e.g.\ A5 in one district and $R$ affinely independent points in
every other.
\end{result}
\begin{plain}
District baselines are identified from each district's own regression. Under M1
the curvature of the regression must be the same in every district: that is the
testable part. The shared contextual split $K$ is not testable. A district
needs only a few compositions to fix its baseline, but they must not lie on a
line (for three groups).
\end{plain}
\proofref{Appendix~\ref{app:phase2b}, \S2, proofs of (a), (b), (c), (d)}
\begin{revnote}
(R-19, R-30, P-B39.) Part (d) and Remark~D1.1 were added: ``$R$ distinct
compositions'' is not enough, they must be affinely independent. The number of
overidentifying restrictions is $(D^\star-1)(C-1)R(R-1)/2$. A5-d is stated on
the support of the restricted measure.
\end{revnote}

\begin{result}[Lemma D2, the estimand under districts --- PHASE2b \S2]
\status{Proved.} $B^*_{rc}=\mu_r^{-1}\sum_d[\sum_sQ^{(d)}_{c,rs}\mu^{(d)}_{rs}+\sum_{s\ne
r}K^{(d)}_{c,rs}\mu^{(d)}_{rs}]$; under M1 its sensitivity to the common $K_{c,rs}$
is $\mu_{rs}/\mu_r$, as pooled.
\end{result}
\proofref{Appendix~\ref{app:phase2b}, \S2, Lemma~D2 steps 1--4}

\begin{result}[Theorem 1$'$, sharpness for any affine mean model --- PHASE2b \S3]\label{res:t1prime}
\status{Proved} (by reduction to PHASE2a Thm~1; measure-theoretic content not
machine-checked). If $m_\theta(v)$ is Borel (H1), bounded (H2), affine and
uniformly continuous in $\theta$ (H3), and the selection set is unaffected
(H4), then $\theta\in\ThI$ iff $m_\theta$ is valid on $\mathcal V_A$ and
$\langle u,m_\theta(v)\rangle\le\psi(u;v)$ for a.e.\ $v$ and all $u$.
\end{result}
\begin{plain}
The sharpness argument of the pooled model carries over to any model in which
the mean is affine in the unknowns, including all the district rungs.
\end{plain}
\proofref{Appendix~\ref{app:phase2b}, \S3, Theorem~1$'$ (step-by-step correspondence with PHASE2a Thm~1)}
\begin{revnote}(R-20.) Hypotheses H1--H4 listed explicitly.\end{revnote}

\begin{result}[Theorem 2, the rungs --- PHASE2b \S3]\label{res:t2}
\status{Proved.} Under A1, A5-d and $P(X\in A_D)=1$: (a) the sharp set of each rung
is the rung restriction intersected with each district's own sharp condition;
(b) the outer set is a polyhedron, every target bound is an LP, and for M1
$\Theta^{M1}_O=\bigcap_d\Theta^{(d)}_O$ in the common $k$; (c) the ladder
\[
\mathcal B^{L6}\subseteq\mathcal B^{L5}\subseteq\mathcal B^{L4(r)}\subseteq\mathcal B^{M1}\subseteq\mathcal B^{M1(0)}\subseteq\mathcal B^{M1(\delta)}\subseteq\mathcal B^{M1(\delta')}\subseteq\mathcal B^{M2}\subseteq DD^*
\]
holds for $0\le\delta\le\delta'$, sharp and outer, whenever the smaller rung's
model is non-empty.
\end{result}
\begin{plain}
Each stronger assumption gives a bound inside the weaker one's, so the ladder can
be reported as a table and readers can see what each assumption buys. M0 is not
on this chain.
\end{plain}
\begin{sketch}
Each rung is an affine mean model, so Thm~1$'$ applies district by district;
the outer version replaces the sharp condition by vertex non-negativity;
nesting follows because each step restricts the same polyhedral system further.
\end{sketch}
\proofref{Appendix~\ref{app:phase2b}, \S3, proofs of (a) steps 1--2, (b) steps 1--4, (c) steps 1--7}
\begin{revnote}
(R-21.) M1(0) is a \emph{sensitivity benchmark}: it keeps the common skew part
but drops the testable symmetric restrictions of D1(b).
\end{revnote}

\begin{result}[Proposition D3, M1 with two rows --- PHASE2b \S4]\label{res:d3}
\status{Proved.} For $R=2$ and vertices $t\in\{l_d,u_d\}$ with $0<l_d\le u_d<1$,
$\Theta^{M1}_O=\{\ell\le\kappa\le h,\sum_c\kappa_c=0\}$ with
$\ell_c=\max_d\max_t[\lambda_c(t)-\zeta_{c,d,1}/(1-t)]$ and
$h_c=\min_d\min_t[\rho_c(t)+\zeta_{c,d,2}/t]$. Pure-composition vertices give
testable restrictions and are dropped from the max/min as stated.
\end{result}
\begin{plain}
A district raises the lower bound when its baseline for group 1 is low and it has
polls that are almost all group 1: there group 1's rate is nearly observed, and
the shared contextual effect carries that information to every district. One
riding with lopsided polls and an extreme baseline can tighten the province-wide
bound.
\end{plain}
\proofref{Appendix~\ref{app:phase2b}, \S4, derivation steps 1--5; pure-composition case; checked against the LP (126 and 195 instances)}
\begin{revnote}(R-22.) $0<l_d\le u_d<1$ made visible; pure-composition vertices derived.\end{revnote}

\begin{result}[Proposition D4, when dummies matter --- PHASE2b \S4]\label{res:d4}
\status{Proved.} (a) If $\zeta\equiv0$, M1's outer set equals the pooled outer set
on $\conv\bigcup_dA_d$. (b) Under M1, $\E[\beta\mid X,N]=M_H(X)+\E[Z_D\mid X,N]$,
and M0 holds iff $\E[Z_D\mid X,N]$ is a.s.\ a function of $X$ alone, affine on
the support; $D\perp(X,N)$ suffices; $D\perp X$ alone does not (Remark~D4.1).
\end{result}
\begin{plain}
The pooled model is safe when district baselines, given composition and unit
size, vary at most linearly with composition. When strongholds also have
different baselines (a popular local candidate in the strongest ridings), or when
district baselines correlate with unit size, the pooled fit mistakes baseline
differences for contextual effects and can miss the truth. This is an
ecological omitted-variable bias.
\end{plain}
\proofref{Appendix~\ref{app:phase2b}, \S4, (a) steps 1--2, (b) steps 1--6; Remark~D4.1}
\begin{revnote}
(R-23.) The earlier version conditioned on $X$ only and claimed $D\perp X$
suffices; the size-weighted estimand uses $\E[\beta\mid X,N]$. Counterexample
with district correlated with size; check in \texttt{verify\_phase2b.py} [F].
Conditional on the observed design, P3.6(c) tests the stronger condition
$\zeta=0$ (Section~\ref{sec:inference}).
\end{revnote}

\begin{result}[Propositions D5--D7 --- PHASE2b \S4]
\status{Proved.}
\textbf{D5}: under M2 the bound is the Minkowski sum of the district bounds, so
widths add.
\textbf{D6}: under M1($\delta$) the bounds are nested in $\delta$, equal M2 for
$\delta\ge\bar\delta$, have continuous piecewise-linear endpoints, and equal
the M1 outer bound at $\delta=0$ (computed with the same $Q^{(d)},\zeta$).
\textbf{D7}: L5 point-identifies $B^*_{rc}=\sum_d(\mu^{(d)}_r/\mu_r)\alpha_{c,d,r}$
(Goodman within districts); L4($r$) pins $K_{c,rs}=Q_{c,rr}-Q_{c,rs}$ for
$s\ne r$, i.e.\ $(C-1)(R-1)$ coordinates, and point-identifies when $R=2$.
\end{result}
\begin{plain}
District-specific lines (M2) are safe but rarely much tighter than DD. The dial
$\delta$ traces a smooth path between M2 and M1, so a reader can see how much
the shared-effects assumption is doing. The constancy rungs give single numbers;
reported beside the bounds they show what ``no contextual effects'' implies
(Remark~D7.1: in the $3\times3$ example L5 gives 0.870 against a truth of
0.836, inside M1 but wrong).
\end{plain}
\proofref{Appendix~\ref{app:phase2b}, \S4, Propositions~D5 (steps 1--3), D6 (steps 1--4), D7 (a) and (b)}

\paragraph{Remarks 6.1--6.3 (PHASE2b \S5).} Remark~6.1: everything above is
identification; with about ten polls per riding the plug-in intersection is
biased inward (72\% coverage in PHASE1 \S7.1), which Phase~3 corrects.
Remark~6.2: with election periods as districts, M1 is ``contextual effects
unchanged'' and leads to Prop.~D8 (Section~\ref{sec:did}). Remark~6.3: the
grouping need not be geographic.

\paragraph{Worked examples (PHASE2b \S6).} In the $2\times2$ two-riding example
(truth 0.819) DD is $[0.453,0.879]$, M0 is empty (and the package reverts to
DD), M2 is $[0.563,0.870]$ and M1 is $[0.618,0.866]$. In the $3\times3$
six-district example the ladder shrinks monotonically, contains the truth at every
rung, and M2$\to$M1 narrows loyalty from width 0.525 to 0.450.
